% !TeX root = ../../Book2.tex
% 纲要标题：### 18.5 周期与指数
% 纲要要点（供目录核对）：
%
% * 周期整除指数，指数整除次数
% * 十六维中心除代数的周期二、指数四例子
%
% <!-- 短节：不设小节 -->
%
% ---

\section{周期与指数}
\label{b2:sec:1805}

指数是除代数代表的次数，也等于有限分裂域扩张次数的最小值；周期则是 Brauer 类的加法阶。前节的半线性矩阵把两种量放在同一计算中：取矩阵行列式，会使因子集的一个幂成为余边。

\begin{theorem}\label{b2:thm:period-divides-index}
设 \(k\) 为域，\(A\) 为中心单 \(k\) 代数，则其 Brauer 类 \([A]\in\operatorname{Br}(k)\) 具有有限阶，且
\[
\operatorname{per}(A)\mid\operatorname{ind}(A)\mid\deg A.
\]
\end{theorem}
\begin{proof}
取同一 Brauer 类中的除代数 \(D\)，令 \(d=\deg D=\operatorname{ind}(A)\)，再取一个有限 Galois 分裂域 \(K/k\)。由\cref{b2:prop:split-algebra-factor-set}，对 \(D_K\cong M_d(K)\) 选矩阵 \(C_g\)，得到
\[
C_g\cdot g(C_h)=c(g,h)C_{gh},\qquad [D]=-[B_c].
\]
取矩阵行列式，令 \(b_g=\det C_g\)，则
\[
c(g,h)^d=b_g\cdot g(b_h)b_{gh}^{-1}.
\]
因此 \(c^d\) 为余边，在 \(H^2(G,K^\times)\) 中为零类。\cref{b2:thm:relative-brauer-cocycles}给出 \(d[B_c]=0\)，所以 \(d[A]=d[D]=0\)。元素的最小正阶于是存在，且整除任何零化它的正整数，特别整除 \(d\)。若 \(A\cong M_r(D)\)，则 \(\deg A=rd\)，还得到第二个整除关系。证明中没有除以 \(d\)，故包括正特征。
\end{proof}

\begin{example}[周期为二、指数为四]\label{b2:exm:period-two-index-four}
在域 \(k=\mathbb R((x))((y))\) 上，代数
\[
D=\mathbb H\otimes_{\mathbb R}k\ \otimes_k\ (x,y)_k
\]
为除代数，且 \(\operatorname{ind}(D)=4\)、\(\operatorname{per}(D)=2\)。这里 \((x,y)_k\) 为四元数代数。
\end{example}
\begin{proof}
给出一个直接的除环模型。先取中心不定元 \(u\)，组成 Laurent 级数除环 \(F=\mathbb H((u))\)。非零级数可以提出其首项，再将剩余的 \(1+h\)、\(h\) 只有正次数，按 \(1-h+h^2-\cdots\) 求逆；每个固定次数的系数只涉及有限多项，所以这个形式逆有意义。

令 \(\tau\) 固定 \(\mathbb H\)，将 \(u\) 送到 \(-u\)。在所有下端有界的形式级数 \(\sum_m a_m v^m\)、\(a_m\in F\) 上规定
\[
(a v^r)(b v^s)=a\tau^r(b)v^{r+s}.
\]
固定次数的系数仍只涉及有限和；\(\tau\) 是自同构，所以三项乘积的两种括号给出相同系数。这便定义结合环 \(E=F((v;\tau))\)。非零首项 \(a_m v^m\) 可逆，其逆为 \(\tau^{-m}(a_m^{-1})v^{-m}\)；提出首项后，再用同样的正次数几何级数构造逆。因此 \(E\) 是除环。

在 \(E\) 中令 \(x=u^2,y=v^2\)。二者都在中心，且 \(F\) 中按 \(u\) 的奇偶次数及四元数基逐项分组，得到左 \(\mathbb R((x))\) 空间的直和
\[
F=\bigoplus_{\epsilon=0}^{1}
\ \bigoplus_{h\in\{1,\mathbf i,\mathbf j,\mathbf k\}}
\mathbb R((x))\,h u^\epsilon.
\]
由于 \(\tau^2=\operatorname{id}\)，\(y=v^2\) 与 \(F\) 交换。再按 \(v\) 的奇偶次数分组，就有
\[
E=F((y))\oplus F((y))v
=\bigoplus_{\delta=0}^{1}
\ \bigoplus_{\epsilon=0}^{1}
\ \bigoplus_{h\in\{1,\mathbf i,\mathbf j,\mathbf k\}}
k\,h u^\epsilon v^\delta,
\qquad k=\mathbb R((x))((y)).
\]
两次形式 Laurent 展开的系数唯一，且 \(1,\mathbf i,\mathbf j,\mathbf k\) 在 \(\mathbb H\) 中线性无关，因此所列十六个元素确为 \(E\) 的 \(k\) 基。

中心也可逐项确定。若 \(z=\sum_m a_m v^m\) 与全部四元数常数交换，因 \(v\) 与这些常数交换，Laurent 展开的唯一性迫使每个 \(a_m\) 与 \(\mathbb H\) 交换。再逐项展开 \(a_m\) 的 \(u\) 系数，得到 \(a_m\in\mathbb R((u))\)。若 \(z\) 还与 \(u\) 交换，由 \(v^m u=(-1)^muv^m\) 及特征为零，所有奇数 \(m\) 的系数都为零。最后与 \(v\) 交换，使每个余下的 \(a_{2r}\) 被 \(\tau\) 固定；其 \(u\) 展开只能含偶数次幂，故 \(a_{2r}\in\mathbb R((x))\)。因此 \(Z(E)=\mathbb R((x))((y))=k\)。

四元数常数与 \(u,v\) 交换，而 \(u^2=x,v^2=y,vu=-uv\)。这些关系及所列基给出 \(E\cong\mathbb H_k\otimes_k(x,y)_k=D\)。因此 \(D\) 自身为十六维中心除代数，指数为四。由\eqref{b2:eq:quaternion-conjugation}，每个四元数代数都通过共轭同构于自身的反代数，其 Brauer 类的二倍为零，所以 \(2[D]=0\)。\(D\) 是不等于 \(k\) 的除代数，Brauer 类非零，故周期恰为二。
\end{proof}
